\documentclass[12pt, letterpaper]{article}

\usepackage[utf8]{inputenc}
\usepackage[T1]{fontenc}
\usepackage[spanish]{babel}
\usepackage{geometry}
\usepackage{amsmath, amssymb, amsthm}
\usepackage{booktabs}
\usepackage{array}
\usepackage{microtype}
\usepackage{setspace}
\usepackage{parskip}
\usepackage{titlesec}
\usepackage{fancyhdr}
\usepackage{enumitem}
\usepackage{bm}
\usepackage{mathtools}
\usepackage{tcolorbox}
\usepackage{hyperref}

\tcbuselibrary{skins, breakable}

\geometry{top=2.8cm, bottom=3.0cm, left=3.2cm, right=3.2cm, headheight=14pt}
\setstretch{1.20}

\pagestyle{fancy}
\fancyhf{}
\fancyhead[C]{\small\textsc{Econometria · Gujarati --- Ejercicio 12.3}}
\fancyfoot[C]{\small\thepage}
\renewcommand{\headrulewidth}{0.4pt}
\renewcommand{\footrulewidth}{0pt}

\titleformat{\section}
  {\large\bfseries\scshape}{\thesection.}{0.6em}{}
  [\vspace{-4pt}\rule{\linewidth}{0.4pt}]
\titleformat{\subsection}{\normalsize\bfseries}{\thesubsection.}{0.5em}{}
\titleformat{\subsubsection}{\normalsize\itshape}{}{0em}{}

\newtcolorbox{clave}{
  colback=white, colframe=black,
  boxrule=0.5pt, arc=3pt,
  left=8pt, right=8pt, top=6pt, bottom=6pt,
  breakable
}

\newcommand{\Var}{\operatorname{Var}}

\begin{document}

\begin{center}
  {\LARGE\bfseries Participacion Laboral en la Industria Metalurgica:}\\[6pt]
  {\LARGE\bfseries Autocorrelacion Pura vs.\ Sesgo de Especificacion}\\[12pt]
  {\large\itshape Modelos de tendencia lineal y cuadratica, estadistico de Durbin-Watson}\\
  {\large\itshape y diagnostico de la fuente de correlacion serial}\\[14pt]
  \rule{0.55\linewidth}{0.6pt}\\[8pt]
  {\normalsize Ejercicio~12.3 --- \textit{Econometria}, Gujarati \& Porter, 5.a ed.\
  por Emanuel Quintana Silva}\\[4pt]
  {\small\textsc{Basado en Gujarati (1969), ``Labor's Share in Manufacturing Industries'',
  \textit{ILRR}, vol.~23, num.~1}}
\end{center}

\vspace{10pt}

% ─────────────────────────────────────────────────────────────────────────────
\section{Planteamiento del problema}
% ─────────────────────────────────────────────────────────────────────────────

Se estudia el movimiento en la participacion laboral ($Y_t$) en el valor agregado
de la industria metalurgica basica, con datos anuales de 1949 a 1964 ($n = 16$
observaciones). Se comparan dos especificaciones:

\begin{align}
  \textbf{Modelo A:} \quad
    Y_t &= \beta_0 + \beta_1 t + u_t
  \label{eq:modA}\\[6pt]
  \textbf{Modelo B:} \quad
    Y_t &= \alpha_0 + \alpha_1 t + \alpha_2 t^2 + u_t
  \label{eq:modB}
\end{align}

Los resultados estimados son:

\begin{align}
  \hat Y_t &= 0.4529 - 0.0041\,t
  \tag{A}\label{eq:estA}\\
  &\quad\quad (-3.9608) \qquad R^2 = 0.5284 \quad d = 0.8252 \notag\\[8pt]
  \hat Y_t &= 0.4786 - 0.0127\,t + 0.0005\,t^2
  \tag{B}\label{eq:estB}\\
  &\quad\quad (-3.2724)\quad (2.7777) \qquad R^2 = 0.6629 \quad d = 1.82 \notag
\end{align}

donde las cifras entre parentesis son las razones $t$.

% ─────────────────────────────────────────────────────────────────────────────
\section{Inciso (a): Correlacion serial en los modelos A y B}
% ─────────────────────────────────────────────────────────────────────────────

Para $n = 16$ observaciones, los valores criticos de la tabla de Durbin-Watson
al nivel de significancia del $5\%$ son:

\begin{center}
\renewcommand{\arraystretch}{1.3}
\small
\begin{tabular}{cccc}
\toprule
\textbf{Modelo} & $k'$ & $d_L$ & $d_U$ \\
\midrule
A (tendencia lineal)    & 1 & 1.106 & 1.371 \\
B (tendencia cuadratica)& 2 & 0.982 & 1.539 \\
\bottomrule
\end{tabular}
\end{center}

\subsection{Modelo A}

\textbf{Comparacion:}
\begin{equation}
  d_A = 0.8252 \;<\; d_L = 1.106
\end{equation}

El estadistico $d$ cae por debajo del limite inferior. Se rechaza la hipotesis
nula de no autocorrelacion.

\textbf{Conclusion:} El Modelo A presenta \textbf{autocorrelacion positiva}
de primer orden estadisticamente significativa al nivel del $5\%$.

\begin{clave}
  Modelo A: $d = 0.8252 < d_L = 1.106$ $\Rightarrow$
  Rechazar $H_0$. Autocorrelacion positiva significativa.
\end{clave}

\subsection{Modelo B}

\textbf{Comparacion:}
\begin{equation}
  d_B = 1.82 \;>\; d_U = 1.539
\end{equation}

El estadistico $d$ supera el limite superior y se encuentra en la region de no
rechazo ($d_U < d < 4 - d_U$, puesto que $4 - 1.539 = 2.461$ y
$1.539 < 1.82 < 2.461$). No se rechaza la hipotesis nula.

\textbf{Conclusion:} El Modelo B \textbf{no presenta evidencia} de correlacion
serial al nivel del $5\%$. El estadistico $d = 1.82$ es consistente con la
ausencia de autocorrelacion.

\begin{clave}
  Modelo B: $d_U = 1.539 < d = 1.82 < 2.461 = 4 - d_U$
  $\Rightarrow$ No rechazar $H_0$. Sin evidencia de autocorrelacion.
\end{clave}

% ─────────────────────────────────────────────────────────────────────────────
\section{Inciso (b): Explicacion de la correlacion serial en el Modelo A}
% ─────────────────────────────────────────────────────────────────────────────

La correlacion serial detectada en el Modelo A es, muy probablemente,
\textbf{autocorrelacion inducida por un error de especificacion} en la forma
funcional, no autocorrelacion pura.

\subsection{Argumento estadistico}

El termino cuadratico $t^2$ en el Modelo B tiene una razon $t = 2.7777$,
estadisticamente significativa al nivel del $5\%$ (valor critico $\approx 2.16$
para $n - k = 13$ grados de libertad). Esto indica que la participacion laboral
sigue una trayectoria \emph{no lineal} en el tiempo y que el Modelo A omite
ese componente de curvatura.

\subsection{Mecanismo de induccion}

Cuando la verdadera forma funcional es cuadratica pero se estima un modelo
lineal, el residuo del Modelo A absorbe el patron no lineal omitido:
\begin{equation}
  \hat u_t^{(A)} = (Y_t - \hat Y_t^{(A)})
  \approx \alpha_2 t^2 + \text{error puro}
\end{equation}
Dado que $\alpha_2 t^2$ es una funcion suave y monotona en $t$, los residuos
de periodos consecutivos seran sistematicamente similares en signo y magnitud,
lo que genera una correlacion serial aparente. En terminos geometricos: si los
datos siguen una curva y el modelo ajusta una recta, los residuos estaran
primero todos positivos, luego negativos, luego positivos de nuevo, produciendo
exactamente el patron que detecta la prueba de Durbin-Watson.

\subsection{Evidencia empirica}

El siguiente cuadro resume la mejora al pasar del Modelo A al Modelo B:

\begin{center}
\renewcommand{\arraystretch}{1.3}
\small
\begin{tabular}{lccc}
\toprule
\textbf{Modelo} & $R^2$ & $d$ & \textbf{Diagnostico} \\
\midrule
A (lineal)    & 0.5284 & 0.8252 & Autocorrelacion positiva \\
B (cuadratico)& 0.6629 & 1.82   & Sin autocorrelacion \\
\midrule
\multicolumn{2}{l}{Mejora al agregar $t^2$}
  & $\Delta d = +0.9948$ & Desaparece la correlacion serial \\
\bottomrule
\end{tabular}
\end{center}

El aumento de $d$ de $0.8252$ a $1.82$ al incluir $t^2$ es la evidencia
mas directa de que la correlacion serial del Modelo A era consecuencia de la
forma funcional incorrecta, no de un proceso autorregresivo en el error.

\begin{clave}
  La participacion laboral sigue una trayectoria cuadratica. Al omitir $t^2$
  en el Modelo A, los residuos absorben la curvatura y exhiben un patron
  serial sistematico. Corregir la especificacion elimina la autocorrelacion.
\end{clave}

% ─────────────────────────────────────────────────────────────────────────────
\section{Inciso (c): Distincion entre autocorrelacion pura y sesgo de especificacion}
% ─────────────────────────────────────────────────────────────────────────────

Distinguir entre autocorrelacion \emph{pura} (inherente al proceso generador
de los errores) y autocorrelacion \emph{inducida} (por mala especificacion del
modelo) es una de las tareas mas delicadas del diagnostico econometrico. Los
criterios siguientes son los mas utilizados en la practica.

\subsection{Criterio 1: respuesta del estadistico $d$ ante cambios de especificacion}

Si al corregir la forma funcional o incluir variables omitidas el estadistico
$d$ pasa de ser significativo a un valor cercano a 2 (sin autocorrelacion), la
fuente era \textbf{sesgo de especificacion}.

Si, por el contrario, $d$ permanece bajo tras agotar las especificaciones
teoricamente justificadas, es mas probable que se trate de \textbf{autocorrelacion
pura}. Este es precisamente el caso del ejercicio: $d$ pasa de $0.8252$ a $1.82$
al corregir la especificacion, lo que senala mala especificacion como causa.

\subsection{Criterio 2: analisis grafico de residuos}

Una inspeccion visual de los residuos contra el tiempo o contra $\hat Y_t$ puede
revelar el origen del problema:
\begin{itemize}[leftmargin=1.6em, itemsep=3pt]
  \item \textbf{Patron en forma de U o de S}: sugiere forma funcional incorrecta
    (componente polinomial omitido). Tipico de sesgo de especificacion.
  \item \textbf{Oscilaciones suaves y regulares sin patron en U/S}: mas consistente
    con autocorrelacion pura de primer orden ($u_t = \rho u_{t-1} + \varepsilon_t$).
  \item \textbf{Patron ciclico de largo plazo}: puede indicar autocorrelacion de
    orden superior o un componente estacional omitido.
\end{itemize}

\subsection{Criterio 3: prueba de Ramsey RESET}

La prueba RESET (Regression Equation Specification Error Test) contrasta
formalmente si potencias de $\hat Y_t$ tienen poder explicativo adicional:
\begin{equation}
  H_0: \text{la forma funcional es correcta}
  \qquad \text{vs.} \qquad
  H_1: \text{existe mala especificacion}
\end{equation}
Un rechazo de $H_0$ senala sesgo de especificacion antes de atribuir la
correlacion serial a autocorrelacion pura.

\subsection{Criterio 4: prueba de Breusch-Godfrey}

La prueba LM de Breusch-Godfrey contrasta autocorrelacion de orden $p$ sin las
restricciones de la prueba de Durbin-Watson. Si tras corregir la especificacion
la prueba LM aun rechaza $H_0$, la autocorrelacion residual es de naturaleza pura.
Si no rechaza, confirma que el problema original era de especificacion.

\subsection{Aplicacion al ejercicio}

En el presente caso, la combinacion de los criterios 1 y 2 es concluyente:
\begin{enumerate}[leftmargin=1.8em, itemsep=4pt]
  \item El estadistico $d$ mejora drasticamente ($0.82 \to 1.82$) al agregar $t^2$.
  \item El coeficiente de $t^2$ es altamente significativo ($t = 2.78$).
  \item El $R^2$ aumenta de $0.53$ a $0.66$, confirmando mejor ajuste.
\end{enumerate}
Estos tres hechos conjuntos indican que la autocorrelacion del Modelo A era
\textbf{inducida por la omision de la tendencia cuadratica}, no un fenomeno
estructural del proceso de error.

\begin{clave}
  \textbf{Regla practica:} si agregar variables o corregir la forma funcional
  elimina la correlacion serial, el problema era de especificacion.
  Si la correlacion persiste tras la correccion, es autocorrelacion pura
  y se debe tratar con GLS, FGLS o errores estandar de Newey-West.
\end{clave}

\vspace{14pt}
\noindent\rule{\linewidth}{0.4pt}

\noindent{\small\textit{Nota.}
El ejercicio ilustra un principio metodologico fundamental: ante un estadistico
$d$ significativo, la primera accion no debe ser corregir la autocorrelacion
sino \emph{revisar la especificacion del modelo}. Aplicar GLS o transformar
las variables cuando el verdadero problema es la omision de un regresor produce
estimadores inconsistentes y no resuelve el problema de fondo.
La jerarquia correcta es: (1) diagnosticar la especificacion, (2) corregir si
procede, (3) solo entonces tratar la autocorrelacion si persiste.}

\end{document}
